Self-improving agents are state-of-the-art (SOTA) on agentic coding benchmarks and have recently been extended to general domains. However, their search methods generally assume a stationary evaluation criterion: a fixed verifier, benchmark, or labeled dataset that remains valid as the agent improves. This ignores a central feature of evolution: species do not optimize against a static environment, but adapt as their environments change with them. We aim to bring the same principle to recursive self-improvement, making evaluation part of the improvement loop and opening search to evolving evaluators, adversarial objectives, and dynamic utilities that may surpass static benchmarks.
We introduce the \textbf{Red Queen G\"odel Machine (\ourmethod)}\footnotemark{}, an evolutionary framework for recursive self-improvement under non-stationary utilities. We make this possible through \emph{controlled utility evolution}: search is organized into epochs with a fixed within-epoch evaluation criterion, while the utility can be updated at epoch boundaries, so self-improvement guarantees hold per epoch as the objective evolves across them. This mechanism is instantiated by \emph{co-evolving} the utility itself, using learned evaluators that improve alongside the task agents they guide. The current preprint provides a preliminary empirical investigation. We begin by showing that even on verifiable coding tasks, \theourmethod improves test pass rate over the prior SOTA by adding a complementary agent-as-a-judge code-review signal. Because the reviewer is queried once while standard coding agents require multi-turn evaluation, this signal is cheaper and \theourmethod uses \textbf{1.35}$\boldsymbol{\times}$--\textbf{1.72}$\boldsymbol{\times}$ fewer tokens. We then turn to scientific paper writing and reviewing, and Olympiad-level proof writing and grading, where \theourmethod improves performance over prior self-improving agents: co-evolved writers reach \textbf{1.78}$\boldsymbol{\times}$--\textbf{1.86}$\boldsymbol{\times}$ higher acceptance rates under a diverse agent-as-a-judge panel, while co-evolved graders reach \textbf{9}$\boldsymbol{\%}$ higher ground-truth accuracy. Because \theourmethod can modify the search objective across epochs, it can also regularize the search process itself. In paper reviewing, the strongest baseline reviewer over-accepts \AI-generated papers at up to \textbf{1.91}$\boldsymbol{\times}$ the human rate. We correct this with \theourmethod by introducing an adversarial objective that discovers reviewers equally stringent on \AI{} and human work. We intend to explore longer search horizons in future updates.